\chapter{Related Work}
\label{ch:related}

\section{Adaptive RAG}
\label{sec:adaptive-rag}

Retrieval-Augmented Generation (RAG)~\cite{rag_survey} augments a parametric language model with a non-parametric retriever, allowing the model to condition its generation on external evidence rather than relying only on parameters learned during pretraining. A typical RAG pipeline can be described as five stages:

\begin{enumerate}
    \item \textbf{Indexing.} The source corpus is segmented into passages, and each passage is stored in a form that supports retrieval, such as sparse lexical indexes or dense vector representations.
    \item \textbf{Query manipulation.} The input question may be rewritten, expanded, normalised, or decomposed so that the retriever can better match it with relevant evidence.
    \item \textbf{Retrieval.} The system selects candidate passages from the indexed corpus, commonly using sparse lexical matching, dense vector similarity, or a hybrid retrieval strategy.
    \item \textbf{Post-retrieval processing.} The initially retrieved passages may be filtered, reranked, compressed, or otherwise refined before they are passed to the generator.
    \item \textbf{Generation.} The language model conditions on the selected evidence and produces the final answer, ideally grounding its output in the retrieved passages.
\end{enumerate}

This retrieve-then-generate paradigm is now widely used to make language model outputs more factual in question answering, explanation writing, and other knowledge-intensive tasks. Its central advantage is that factual evidence can be updated externally without retraining the language model. For multi-hop question answering, however, the quality of the retrieved evidence depends strongly on the query that is issued at each step: a single original question may not contain the intermediate entity needed to retrieve the next supporting passage.

This fixed retrieval assumption creates two well-documented failure modes:

\begin{itemize}
    \item \textbf{Noise injection on easy queries.} For questions the model already knows (e.g.\ ``Who wrote \emph{Hamlet}?''), retrieved passages introduce distractor entities that can override correct parametric knowledge~\cite{hallucination_survey}.
    \item \textbf{Insufficient evidence on hard queries.} For multi-hop questions, a single retrieval pass rarely surfaces all bridging entities, leading to incomplete or compositionally wrong answers~\cite{ircot}.
\end{itemize}

These two failure modes motivate adaptive retrieval: deciding \textbf{when} retrieval is \mbox{useful}, \textbf{how much} evidence to fetch, and what signal should determine the gate.

Adaptive RAG methods decide whether and how much to retrieve for each query or reasoning step. They can be characterised along three axes: decision granularity, gating signal, and training cost. Decision granularity ranges from question-level routing to sentence-, token-, or sub-question-level decisions, while training cost ranges from training-free methods to supervised fine-tuning, reinforcement learning, or preference optimisation. Since the gating signal is the most consequential axis for SEEK, we organise existing methods into three families: classifier-based gating, confidence-based gating, and reasoning-step decision methods.

\paragraph{Classifier-based gating.}

The first family trains a dedicated discriminator that inspects the query and emits a retrieval decision.

\textbf{Adaptive-RAG}~\cite{adaptiverag} routes each incoming query to one of three retrieval strategies according to its predicted complexity. Easy questions are answered without retrieval, medium-complexity questions use a single retrieval step, and complex questions are routed to a multi-step retrieve-and-reason pipeline such as \textbf{IRCoT}~\cite{ircot}. The router is a T5-Large classifier trained with automatically derived labels. Each training question is assigned to the cheapest strategy that still produces the correct answer, so the classifier learns a cost-aware mapping from question complexity to retrieval depth. This design reduces unnecessary retrieval compared with always using a multi-step pipeline. However, the decision is made only once at the beginning of the query. As a result, the system cannot react to uncertainty that appears during intermediate reasoning, and its routing labels remain coarse: questions with different numbers of reasoning hops may still be assigned to the same multi-step category.

\textbf{SKR}~\cite{skr} approaches adaptive retrieval from the perspective of self-knowledge. Instead of predicting general question complexity, SKR estimates whether the language model is likely to know the answer from its own parametric memory. If the question is judged as known, the model answers directly; if it is judged as unknown, external retrieval is invoked. This judgement can be implemented with a trained classifier or with memory-based self-knowledge strategies. In the classifier setting, SKR is a binary retrieval gate: it decides whether to retrieve, but does not distinguish between single-step and multi-step retrieval. Therefore, compared with Adaptive-RAG, SKR gives a more direct estimate of whether retrieval is needed, but it still commits to a question-level decision before reasoning begins.


\paragraph{Confidence- and uncertainty-based gating.}

The second family replaces an external router with uncertainty signals observed during generation. These methods are often training-free, since the retrieval decision is derived from token probabilities, entropy, \mbox{attention} patterns, or internal model states rather than from an additional trained classifier.

\textbf{FLARE}~\cite{flare} moves the retrieval decision from the question level to the generation level. At each step, it first drafts a tentative next sentence and checks whether the sentence contains low-confidence tokens. If the confidence falls below a threshold, the tentative sentence is used to formulate a retrieval query, and the model regenerates the sentence with the retrieved evidence. This design allows retrieval to occur only when the model appears uncertain during generation. However, token probability is only an indirect proxy for factual uncertainty. A low-probability token may reflect lexical variation rather than missing knowledge, and sentence-level regeneration may not align well with multi-hop reasoning, where the natural decision unit is often a sub-question.

\textbf{DRAGIN}~\cite{dragin} further refines dynamic retrieval by considering both when to retrieve and what to retrieve. Its \emph{Real-time Information Need Detection}(RIND) module estimates whether the current generation state requires external evidence, while its \emph{Query Formulation based on Self-attention}(QFS) module constructs retrieval queries from salient \mbox{context} tokens instead of relying only on the most recent sentence. Compared with FLARE, DRAGIN provides a more context-aware retrieval trigger and query formulation strategy. Nevertheless, its decision is still made at the token-generation level, which can be mismatched with explicit sub-question reasoning, and it introduces additional overhead from attention-based query construction.

\textbf{SeaKR}~\cite{seakr} estimates retrieval need from the model's internal states rather than from output probabilities alone. It activates retrieval when the model exhibits high self-aware uncertainty and reranks retrieved passages according to how much they reduce that uncertainty. This makes the retrieval signal more closely tied to the model's internal generation dynamics than token-level probability alone. However, the signal is still an estimate of \mbox{uncertainty} rather than a direct measure of answer correctness. A model may be uncertain for reasons unrelated to missing external knowledge, and retrieval may not necessarily improve the answer even when uncertainty is high.

\textbf{TARG}~\cite{targ} studies adaptive retrieval as a training-free, single-shot decision. It generates a short no-context draft and computes uncertainty from the prefix logits to decide whether retrieval should be invoked. In addition to entropy, TARG considers a top-1/top-2 logit-margin signal and a small-sample variance signal. Its main strength is that it \mbox{provides} a lightweight retrieval gate without training an auxiliary classifier or modifying the language model. However, its evaluation mainly focuses on open-domain factoid QA settings, where each question is typically answered with a single retrieval decision. Its behaviour in multi-hop QA or decomposition-based retrieval settings therefore remains less explored.

Across this family, uncertainty provides a useful and interpretable signal for adaptive retrieval, but it remains an imperfect proxy. Low uncertainty does not guarantee a correct answer, and high uncertainty does not always indicate that retrieval will help. SEEK addresses this limitation by using entropy as a sub-question retrieval signal and as one component of a late verifier, rather than treating it as a single hard router for the entire question.

\paragraph{Learned retrieval-policy methods.}

The third family does not rely on a fixed classifier or a hand-set uncertainty threshold. Instead, it learns retrieval decisions as part of the model's reasoning or generation process.

\textbf{Self-RAG}~\cite{selfrag} trains the language model to generate special reflection tokens that control retrieval and evaluate retrieved passages. For example, the model can emit \texttt{[Retrieve]} to decide whether external evidence is needed, and then use additional tokens to judge whether a passage is relevant, supportive, or useful. This design makes retrieval part of generation itself, rather than a separate preprocessing step. However, it requires fine-tuning on reflection-token supervision and a custom decoding procedure that interleaves normal text generation with retrieval-related decisions.

\textbf{DeepRAG}~\cite{deeprag} formulates retrieval-augmented reasoning as a step-by-step decision process. At each sub-query, the model chooses between answering from parametric memory and issuing a retrieval call. Its training procedure first uses binary tree search to explore both choices for each sub-query, and then applies preference-based calibration to learn which reasoning traces should be preferred. This gives DeepRAG strong sub-question-level retrieval control, but the cost is substantial. Tree-search rollout becomes expensive as the number of reasoning steps grows, preference pairs must be carefully constructed, and the final retrieval policy is learned implicitly rather than exposed as an interpretable score.

A related line of work studies retrieval cost more directly. \textbf{Fast-or-Better}~\cite{fastorbetter} treats the number of retrieved passages or retrieval calls as a controllable trade-off between latency and accuracy, rather than assuming one fixed operating point. This view reinforces an important point for adaptive RAG: retrieval should be evaluated not only by accuracy, but also by the cost required to obtain that accuracy.

SEEK is closest in spirit to DeepRAG because both methods make retrieval decisions at the sub-question level. The key difference is that SEEK avoids tree search and preference optimisation. It uses a single supervised decomposer and replaces the learned retrieval policy with a training-free verifier based on entropy, decomposition stability, answer type, and candidate agreement.

Table~\ref{tab:related-work} summarises the key differences between SEEK and existing adaptive RAG methods along the three axes introduced above. SEEK occupies the under-explored region of \emph{late verifier-based routing with interpretable, training-free signals}: only a single SFT pass is needed for the decomposer, and the final route is decided by transparent query-level diagnostics rather than a learned policy.

\begin{table}[h]
    \centering
    \caption[Comparison of Adaptive RAG Methods]{Comparison of adaptive RAG methods along granularity, gating signal, interpretability, training cost, and decision unit.}
    \label{tab:related-work}
    \setlength{\tabcolsep}{6pt}
\begin{tabular}{ l c c c c c }
        \toprule
        Method & Decomposer & Gating Signal & Interpretable & Training Cost & Granularity \\
        \midrule
        Adaptive-RAG~\cite{adaptiverag} & None       & Query classifier            & Partial & Low (T5-Large) & Question \\
        FLARE~\cite{flare}              & None       & Token probability           & Yes     & None           & Sentence \\
        Self-RAG~\cite{selfrag}         & None       & Reflection tokens           & Partial & High (distill) & Span \\
        DRAGIN~\cite{dragin}            & None       & Entropy + attention         & Yes     & None           & Token \\
        SKR~\cite{skr}                  & None       & Binary self-knowledge       & Partial & Low            & Question \\
        SeaKR~\cite{seakr}              & None       & Hidden-state variance       & No      & None           & Sentence \\
        TARG~\cite{targ}                & None       & Prefix uncertainty          & Yes     & None           & Question \\
        DeepRAG~\cite{deeprag}          & SFT        & Learned policy (SFT+DPO)    & No      & High           & Sub-Q \\
        \textbf{SEEK (Ours)}            & SFT        & Entropy-aware verifier       & \textbf{Yes} & \textbf{Low} (SFT only) & \textbf{Sub-Q + query} \\
        \bottomrule
    \end{tabular}
\end{table}

\section{Reasoning and Knowledge Boundaries in RAG}
\label{sec:reasoning-rag}
\label{sec:knowledge-boundary}

A parallel thread focuses on \emph{how} to retrieve and reason once retrieval is deemed necessary, independent of the gate.

The foundational technique is \textbf{IRCoT}~\cite{ircot}, which interleaves chain-of-thought reasoning with retrieval: after each generated CoT sentence, the partial reasoning becomes the next query. IRCoT is the canonical “multi-step” implementation referenced by \textbf{Adaptive-RAG}~\cite{adaptiverag} and \textbf{DeepRAG}~\cite{deeprag}, but its retrieval count grows linearly in the CoT length, regardless of whether intermediate steps are factually difficult. Question decomposition provides a complementary way to structure multi-hop QA. Rather than retrieving after each free-form reasoning sentence, decomposition-based methods explicitly split a complex question into simpler sub-questions that can be answered independently or sequentially. Early systems such as \textbf{DecompRC}~\cite{decomprc}  show that compositional questions can be reduced to single-hop reading-comprehension steps, while \textbf{BREAK/QDMR}~\cite{break_qdmr} formalises decomposition as an ordered sequence of natural-language operations required to answer the original question. SEEK inherits this decomposition paradigm, but uses it primarily as an inspectable reasoning path: it preserves the original-question context for decomposed steps, evaluates whether the generated sub-questions are useful, and verifies whether the direct or decomposed answer should be trusted.

More recent work synthesises reasoning data to internalise the retrieve-and-reason loop. \textbf{Auto-RAG}~\cite{autorag} autonomously synthesises multi-turn retrieval--reasoning trajectories and fine-tunes on them, so that the model plans its own retrieval steps without hand-crafted prompts. \textbf{Search-o1}~\cite{searcho1} embeds retrieval inside a large reasoning model in an agentic loop, augmenting o1-style long-chain reasoning with on-demand search and a document-reasoning module that distils retrieved passages before they re-enter the chain. \textbf{AirRAG}~\cite{airrag} activates the model's intrinsic reasoning by combining Monte Carlo Tree Search with self-consistency to explore alternative solution paths. These methods substantially strengthen the reasoning side of RAG, but --- like DeepRAG --- they rely on heavy data synthesis or search-time expansion; SEEK instead keeps the reasoning component a single lightweight SFT decomposer and concentrates its novelty in training-free verification.

Every adaptive method above ultimately rests on a model's ability to recognise its own \emph{knowledge boundary} --- the line between what it reliably knows and what it must look up. A dedicated line of work studies this capacity directly, and we group it into three strands: \emph{self-knowledge elicitation}, \emph{uncertainty estimation}, and the \emph{interaction between the knowledge boundary and retrieval}.

\paragraph{Self-knowledge and its calibration.} A first strand asks whether a model can introspect on what it knows. Kadavath et al.~\cite{kadavath2022} show that LLMs are, to a useful degree, calibrated: a model can be prompted to predict the probability \texttt{P(True)} that its own answer is correct, and these self-judgements correlate with actual accuracy --- though calibration degrades on out-of-distribution and adversarial questions. Yin et al.~\cite{yin2023} construct a ``SelfAware'' benchmark of answerable and unanswerable questions and find that LLMs only partially ``know what they don't know'': self-knowledge can be elicited but lags well behind human performance. From the supervised side, Kapoor et al.~\cite{kapoor2024} argue that models ``must be taught to know what they don't know'', demonstrating that lightweight calibration-tuning or probing substantially sharpens a model's estimate of its own uncertainty. Yin et al.~\cite{yin2024} formalise the notion of a knowledge boundary itself, proposing to benchmark a model by probing the frontier of its reliable knowledge rather than by aggregate accuracy alone.

\paragraph{Uncertainty estimation as a boundary signal.} A second strand asks how to \emph{measure} the boundary at inference time. Beyond retrieval-specific methods, recent work on LLM uncertainty quantification compares several confidence estimators for answer correctness, including token entropy, maximum softmax probability, margin-based scores, verbalised confidence, and response consistency~\cite{maqa}. These studies suggest that no single confidence signal is universally reliable across tasks and models. This motivates evaluating whether entropy is sufficient in SEEK's QA setting, or whether a margin-based signal provides a more discriminative retrieval gate.

Semantic entropy further shows why token-level uncertainty can be limited for natural-language answers: different surface forms may express the same meaning, so uncertainty over token sequences does not always correspond to uncertainty over answers. Kuhn et al.~\cite{semanticuncertainty} introduce \emph{semantic entropy}, which clusters multiple sampled answers by meaning before computing entropy, so that paraphrases of one correct answer do not register as uncertainty; Farquhar et al.~\cite{semanticentropy} scale this idea and show in \emph{Nature} that semantic entropy reliably detects confabulations across QA and biography generation. This observation is important for SEEK because exact QA answers often involve aliases, entity variants, and short spans. We therefore treat token entropy as a cheap diagnostic signal rather than as a correctness guarantee.

\paragraph{The boundary under retrieval.} A third strand studies the boundary specifically in the RAG setting. Ren et al.~\cite{ren2023} show that retrieval augmentation not only injects evidence but also \emph{reshapes} the boundary: relevant passages let a model answer questions it could not before, while irrelevant or distracting passages can push a previously-correct answer over the boundary into error --- precisely the ``noise injection on easy queries'' failure mode of \S\ref{sec:adaptive-rag}. This makes a per-step boundary estimate, rather than a one-shot query-level one, essential when retrieval and reasoning are interleaved.

These three strands form the conceptual foundation of SEEK. The failure modes of \S\ref{sec:adaptive-rag} --- spurious retrieval on easy queries, missing retrieval on hard ones, and confident hallucination under a hard entropy threshold --- are all symptoms of a poorly localised knowledge boundary. Where confidence-based methods read a single uncertainty proxy and policy-based methods bury the boundary inside a learned policy, SEEK makes the boundary inspectable through a late verifier: entropy, decomposition stability, and answer agreement jointly decide whether the direct or decomposed path should be trusted.

\section{Evaluation Datasets}
\label{sec:evaluation-datasets}

The evaluation of adaptive RAG hinges on datasets that span both single-hop and multi-hop complexity. This thesis uses HotpotQA, 2WikiMultihopQA, MuSiQue, and PopQA to cover both multi-hop reasoning and single-hop factual recall.

\paragraph{HotpotQA.}

HotpotQA~\cite{hotpotqa} is the canonical multi-hop QA benchmark, containing 113K Wikipedia-based questions requiring 2-hop reasoning. Each example provides (a) the question, (b) the gold answer, (c) two gold-supporting paragraphs out of 10 distractors (in the \texttt{distractor} setting), and (d) sentence-level supporting facts. Questions fall into two compositional patterns: \emph{bridge} (the answer to sub-question 1 is the subject of sub-question 2) and \emph{comparison} (two entities are compared along a shared attribute). HotpotQA is a primary training source for DeepRAG and a standard evaluation set for nearly all adaptive RAG papers.

\paragraph{2WikiMultihopQA.}

2WikiMultihopQA~\cite{twowiki} constructs 192K multi-hop questions by leveraging Wikidata knowledge-graph paths, ensuring each question requires traversal of an explicit reasoning chain. It includes four question types (comparison, inference, compositional, and bridge-comparison) and provides KG-triple annotations alongside textual evidence. Compared to HotpotQA, 2WikiMultihopQA exhibits less surface-form leakage between the question and the supporting passages, making it a more reliable test of compositional reasoning.

\paragraph{MuSiQue.}

MuSiQue (\textbf{Mu}ltihop questions via \textbf{Si}ngle-hop \textbf{Que}stions)~\cite{musique} addresses two pathologies of earlier datasets: (i) shortcut answers via single-hop reasoning, and (ii) weak distractor selection. MuSiQue is built by composing 2--4 single-hop questions into chains using strict connectivity criteria, yielding 25K questions of 2-hop, 3-hop, and 4-hop varieties. Importantly, MuSiQue provides the \emph{decomposed} sub-questions as supervision, which makes it a natural training source for decomposer models such as the one used in SEEK. Its high distractor difficulty makes it the hardest of the four datasets in this thesis.

\paragraph{PopQA.}

PopQA~\cite{popqa} is a \emph{single-hop} entity-centric QA dataset of 14K questions, each derived from a Wikidata subject--property--object triple. Critically, PopQA tags each entity with its \emph{popularity} (measured by monthly Wikipedia page views), enabling stratified evaluation of factual recall on long-tail entities. Despite being single-hop, PopQA exposes the over-decomposition failure mode of decomposer-based pipelines: trained on multi-hop data, a decomposer may emit two-step plans even when one step suffices, harming both EM and efficiency.

\paragraph{Single-hop open-domain QA.}

For completeness, prior adaptive-RAG papers also evaluate on \textbf{Natural Questions (NQ)}~\cite{nq}, \textbf{TriviaQA}~\cite{triviaqa}, and \textbf{SQuAD}~\cite{squad}. These are single-hop datasets whose answers typically exist in a single Wikipedia paragraph, and they primarily test whether an adaptive method can correctly \emph{skip} retrieval rather than orchestrate it. This thesis follows the DeepRAG protocol of evaluating on PopQA as the representative single-hop test set.
